\section{Linear subsequences and heights for isolated components}\label{sec:scale}

At an isolated type, Section~\ref{sec:transfer} has already supplied
sharp boundary measurements $x\mapsto P_x$. We now construct a height
at each source label: a positive injective self-adjoint operator that
commutes with $P_x$ and transforms by the inverse boundary Jacobians.
The construction will take an inverse power of the modulus of a
horizontal derivative of $x\mapsto P_x$. Its main analytic requirement
is a square-integrable derivative on a dense, invariant domain.

The first part of the section controls Fourier lengths. We prove that,
along a sequence of long source intervals, every fixed input retains a
positive amount of mass below a linear length bound. The addition theorem
of Section~\ref{sec:ray} and property~$(T)$ are used to rule out total
superlinear growth. This relatively weak upper-window conclusion is
sufficient: lattice packing costs one factor of the length, which is
canceled by averaging over a source stabilizer interval of comparable
length. The resulting estimate for ordered products at opposite source
labels controls small horizontal differences. We then construct the
closed derivative modulus and prove that its kernel is zero.

The source ray belongs to an entire simple factor. The isolated factors
are real and locally isomorphic to $\Sp(n,1)$, $n\ge2$, or
$F_{4(-20)}$. Complementary factors may be archimedean or nonarchimedean.
We retain a compact window in those complementary coordinates in each
lattice count. The horizontal metrics and angular probabilities are the
compatible compact-invariant choices made in Section~\ref{sec:ray}.

We use the admissibility and probability-law conventions of
Definition~\ref{def:ray-admissible} and Section~\ref{sec:ray}.  Thus all
varying finite factor traces are normalized, the base $Y=G/\Gamma$ has
probability measure, and extra finite matrices carry the unnormalized
matrix trace with fixed matrix size.  Every application of the ray
estimates below has both the stated operator bounds and bounded
integrated Hilbert norms.  A lower effect bound will always concern a
fixed Hilbert vector; it is not a bound uniform over a moving unit
ball.

Here are the precise earlier geometric conclusions used.  The
complementary regular realizations and all-factor mixing are supplied
by Lemma~\ref{lem:setup-normal-tensors},
Proposition~\ref{prop:setup-induced-algebras}, and
Lemma~\ref{lem:setup-product-mixing}. The gap count in
Proposition~\ref{prop:links-gap-count}
gives full-gap escape on complete vertex-origin lexicographic chains.
The angular conclusions of Propositions~\ref{prop:transfer-types},
\ref{prop:transfer-isolated-sharpness}, and~\ref{prop:transfer-sharp-flags}
give sharp, strongly continuous
pure-vertex measurements, their source and measured covariance,
nonconstancy, a permutation of isolated types, and tightness of every
complementary length on a pure isolated ray. We also use scalarity of
limiting rescaled one-column effects from
Lemma~\ref{lem:links-invariant-variables}: their measured invariance places
them in $\mathcal N$, and their face-parabolic and complementary source
invariances make them scalar by all-factor mixing.  The compact-change
argument for those effects uses bounded approximations of the fixed
transporters and has arbitrary fixed slack in radial windows.  These
are conclusions of the preceding constructions, not additional
assumptions on the lattice or its projected actions.

\subsection{An admissible model for the original module}

\begin{lemma}\label{lem:scale-original-model}
For the two-slot model of Section~\ref{sec:setup} there is a
finite tensor \(D=B\bar\otimes L_\eta(\Gamma)\), with \(B\)
diffuse, and an admissible representation
\(\delta:\Gamma\to\mathcal U(D)\) for each isolated measured
factor, obtained from the source action by a fixed finite
matrix corner equivalence.  It has a sequence \(\gamma_n\)
whose Fourier laws escape in every measured simple factor.
Rescaled continuous radial laws, scalar effects, and
complementary tightness are preserved by this corner equivalence.
\end{lemma}

\begin{proof}
The complementary \(\Lambda\)-factor in the two-slot regular
model is diffuse.  In a sufficiently large finite matrix
amplification choose a projection in that coefficient factor
with the same trace as the regular-model projection \(e\).
A partial isometry identifying the projections gives the
required unital tensor-coordinate representation.  Scalar
relative commutant follows from
Proposition~\ref{prop:setup-induced-algebras}.

Choose a full lexicographic chain starting at a vertex and adding
one new simple gap at each stage. By
Proposition~\ref{prop:links-gap-count}, all its target gaps diverge.
Apply this to the finite
matrix of unit columns, and choose successively finite path
times at which the integrated mass in successively larger
Cartan windows is small.  At each selected time, first make the sum of the errors for the
finite matrix and the finitely many prescribed windows smaller than a
number tending to zero. Choose a base point where that nonnegative
sum is no larger than its integrated bound. Each individual error is
then small at the same base point. The resulting lattice cocycle values
give \(\gamma_n\).

A fixed matrix corner equivalence preserves this escape.
Approximate its entries and their adjoints in trace Hilbert
norm by uniformly bounded finite Fourier matrices.  Errors
are uniformly small on the uncut unitary columns.  Each
fixed Fourier shift changes every Cartan coordinate by a
bounded amount.  It therefore preserves escape, and it
intertwines rescaled continuous radial tests in the limit.
After moving the test through the approximating matrix,
its uncut Gram factor converges to the Gram of the partial
isometry.  The same argument, with slack in an unscaled
window, preserves complementary tightness.

For amplified admissibility, finite-rank compressions of
\(\lambda_{X_T}(h_T)\) tend to zero when \(r_T(h)\to\infty\).
This is also seen by embedding \(L^2(X_T)\) into a multiple
of \(L^2(T)\), since the maximal compact stabilizer is compact,
and testing compactly supported vectors.  Compress
\(\Delta_T(\delta(\gamma_n))\) between finite-rank operators
in the translation coordinate.  Plancherel and the escaping
Fourier law make its tracial \(L^2\)-norm tend to zero.
Pairings for conjugation between finite-rank, bounded
trace-Hilbert tensors consequently tend to zero, by
Cauchy--Schwarz; approximation extends this to all
trace-Hilbert vectors.  An invariant vector would have
constant nonzero self-pairing along this sequence.
\end{proof}

\begin{lemma}[Entropy lower bound]\label{lem:scale-entropy}
Suppose an isolated source type \(i\), with factor \(S\),
is matched to the measured type \(p(i)\), with factor \(T\).
Let \(Q_S\) and \(Q_T\) be the homogeneous dimensions of their boundaries.
For the representation \(\delta\) of
Lemma~\ref{lem:scale-original-model}, complementary lengths
are uniformly tight on the pure source ray, and, for every
\(0<c<Q_S/Q_T\),
\begin{equation}
 P_{\delta,t}(t)([0,c])\longrightarrow0.
 \label{eq:scale-entropy}
\end{equation}
Both assertions hold uniformly under source pre- and
post-translations in compact sets.
\end{lemma}

\begin{proof}
Complementary tightness is the pure-ray conclusion of
Section~\ref{sec:transfer}, transported through the fixed corner
equivalence.  We prove the lower bound before
that equivalence.  Complementary regularity supplies a
nonzero constant frame column \(f\) for
\(\pi_0(\Gamma)\) together with \(W_\Lambda\).  If \(e_j\)
are the finitely many identity columns in the measured
regular model, the frame inequality, applied to \(W_h e_j\),
gives
\begin{equation}
 \sum_{\gamma\in\Gamma}
 \|D_{\{h\}}\pi_0(\gamma)f\|_2^2\le C_f
 \qquad(h\in\Gamma).
 \label{eq:scale-bessel-count}
\end{equation}
Indeed the left side is the sum, over \(\gamma\), \(\lambda\),
and the finite column index \(j\), of the squared coefficients
against \(W_hW_\lambda e_j\).  Summing first over
\((\gamma,\lambda)\) is exactly the complementary frame bound.

Let \(\Omega_t\) be a source Cartan shell of fixed positive
thickness about length \(t\) in \(S\), with all compact
rotations and with complementary coordinates in a fixed
compact neighborhood.  Cartan integration gives
\(|\Omega_t|\asymp e^{Q_S t}\).  Restrict both section values
to a compact set \(K\).  For fixed \(y\) and \(\gamma\),
the times with transporter \(\gamma\) lie in
\[
 K\gamma d_0(y)^{-1},
\]
so their Haar volume is at most \(|K|\).  Integrating
\eqref{eq:scale-bessel-count} therefore bounds the averaged
mass at each measured index by a constant independent of
that index and of \(t\).

For a fixed complementary compact window, the number of
measured lattice indices with \(r_T(h)\le c't\) is at most
\(C e^{Q_Tc't}\).  To see this without assuming a discrete
projection to \(T\), thicken the indices on the right by a
small lattice-disjoint neighborhood in the full ambient
group.  Cartan length changes by a bounded amount and the
complementary window remains compact; Haar integration
then gives the stated bound.  It follows that the restricted
low-length mass averaged over \(\Omega_t\), divided by
\(|\Omega_t|\), is at most
\begin{equation}
 C e^{(Q_Tc'-Q_S)t}.
 \label{eq:scale-volume-ratio}
\end{equation}

Choose \(c<c'<Q_S/Q_T\).  The lost section sets have
arbitrarily small measure, uniformly in the translating
time, and the norm of the uncut translated constant column
is unchanged.  Complementary tightness similarly removes
the restriction to the fixed measured compact window.
Finally Lemma~\ref{lem:links-invariant-variables} makes every
limiting low-length effect scalar, and the compact-change
argument identifies its value uniformly throughout the
thick shell, with slack from \(c\) to \(c'\).  A subsequence
violating \eqref{eq:scale-entropy} would thus have positive
averaged mass in \eqref{eq:scale-volume-ratio}, a
contradiction.  The corner-change argument of
Lemma~\ref{lem:scale-original-model} finishes the proof.
\end{proof}

\subsection{Tightening a Kazhdan window}

For a sequence of cutoffs \(L_n\ge1\), a bounded sequence
\(x_n\in D_n\) is called tight at scale \(L_n\) if
\begin{equation}
 \lim_{C\to\infty}\lim_{n\to\mathcal U}
 \|D_{1_{\{r_T>CL_n\}}}x_n\|_2=0,
 \label{eq:scale-tight-sequence}
\end{equation}
where \(\mathcal U\) is a free ultrafilter.  Ultraproducts in
this Subsection are tracial ultraproducts.  Passage to an
ordinary subsequence will be made after all countably many
requirements have been fixed.

\begin{lemma}\label{lem:scale-tight-factor}
The tight elements form a von Neumann subalgebra \(Q_0\)
of \(\prod_{\mathcal U}D_n\), and \(Q_0\) is a factor.
If each \(\Theta_n\) is relatively irreducible, the
ultraproduct representation \(\Theta\) also has scalar
relative commutant.
\end{lemma}

\begin{proof}
Smooth radial truncations have a uniform operator bound
by Proposition~\ref{prop:ray-radial}.  Thus tight
elements admit Hilbert approximation by bounded sequences
supported in \(r_T\le CL_n\).  The triangle inequalities
\eqref{eq:ray-triangle} prove closure under multiplication;
inversion preserves length and gives closure under adjoint.
The tightness condition is closed in Hilbert norm on every
bounded operator ball.  These facts prove von Neumann
closure.

All coefficient-unitary sequences belong to \(Q_0\).  If
\(z\) is central in \(Q_0\), its representing sequence
commutes asymptotically, uniformly, with those unitaries:
otherwise choose a violating unitary at each index.  The
Hilbert closed convex hull of the coefficient conjugation
orbit gives conditional expectation onto its commutant.
Consequently \(z\) is represented, up to a Hilbert-null
sequence, on the group factor alone.

All fixed group labels also belong to \(Q_0\), since
\(L_n\ge1\).  Conjugation by those labels acts on the
orthogonal complement of the identity Fourier coefficient
without invariant vectors.  In fact the squared moduli of
an invariant vector's coefficients would be constant on
conjugacy classes; ICC forces all nonidentity coefficients
to vanish.  Scalar cocycle phases do not change this
argument.  A fixed Kazhdan pair for \(\Gamma\) gives a
uniform spectral gap for all these conjugation
representations, even as the cocycles vary.  Approximate
commutation therefore forces \(z\) to be scalar.

Apply the same Kazhdan inequality to
\(\Ad\Theta_n\) on \(L^2(D_n)\ominus\mathbb C1\).
Relative irreducibility says exactly that this representation
has no invariant vector.  An ultraproduct element commuting
with every \(\Theta(\gamma)\) must accordingly be scalar.
\end{proof}

\begin{lemma}[Kazhdan tightening]\label{lem:scale-tightening}
Fix a finite Kazhdan set \(\mathcal F\subset\Gamma\).
For every \(\varepsilon>0\) there is
\(\varepsilon_K>0\) with the following property.
Suppose \(\Theta_n\) are admissible and
\begin{equation}
 \|D_{1_{\{r_T>L_n\}}}\Theta_n(\gamma)\|_2
 \le\varepsilon_K
 \quad(\gamma\in\mathcal F).
 \label{eq:scale-kazhdan-window}
\end{equation}
After passage to a subsequence there are unitaries \(u_n\)
with \(\limsup_n\|u_n-1\|_2<\varepsilon\) such that every
fixed element of \(u_n^*\Theta_n(\Gamma)u_n\) is tight at
scale \(L_n\).  Admissibility is preserved.
\end{lemma}

\begin{proof}
The Kazhdan gap first supplies an invariant finite right module close to
the original copy of the tight algebra. Its dimension is close to one,
but we need it to be exactly one in order to obtain a conjugating
unitary. We will exclude any excess dimension by applying the ray lift
to a column orthogonal to the tight algebra; the resulting almost
invariant amplified state contradicts admissibility. We then align the
unitary with the identity.

\emph{A nearby invariant module.}
Put \(D_\mathcal U=\prod_{\mathcal U}D_n\) and use
Lemma~\ref{lem:scale-tight-factor}.  In the Jones basic
construction \(\langle D_\mathcal U,e_{Q_0}\rangle\), with
canonical semifinite trace \cite{Jones1983}, one has
\begin{equation}
 \|\Theta(\gamma)e_{Q_0}\Theta(\gamma)^*-e_{Q_0}\|_2^2
 =2\|\Theta(\gamma)-E_{Q_0}\Theta(\gamma)\|_2^2 .
 \label{eq:scale-jones-movement}
\end{equation}
The right side is \(O(\varepsilon_K^2)\): a smooth cutoff
equal to one on \([0,L_n]\) produces a tight approximation.
Its outer radius may be a fixed larger multiple of \(L_n\),
chosen so that the multiplier scale always exceeds the
threshold in Proposition~\ref{prop:ray-radial}.
This does not increase the error beyond the tail in
\eqref{eq:scale-kazhdan-window}.
The Hilbert-space Kazhdan inequality \cite{BHV2008} makes the projection of \(e_{Q_0}\) to the
invariant Hilbert subspace \(O(\varepsilon_K)\)-close to it.
This projection is the Hilbert limit of convex combinations
of its conjugates, and is therefore a positive contraction
\(z\).  Put \(p_0=1_{[1/2,1]}(z)\).  Spectral rounding gives
\begin{equation}
 \|p_0-e_{Q_0}\|_2\le2\|z-e_{Q_0}\|_2=O(\varepsilon_K).
 \label{eq:scale-rounded-projection}
\end{equation}
For completeness, split the left squared norm into
\(\Tr(e_{Q_0}(1-p_0))+\Tr((1-e_{Q_0})p_0)\);
the two terms are bounded by four times the corresponding
terms with \((1-z)^2\) and \(z^2\).  This proves the
inequality without a commutation assumption.
The projection \(p_0\) commutes with \(\Theta(\Gamma)\), and
its trace \(d'\) is close to one.

Take \(\varepsilon_K\) small enough that \(d'<2\).
A right \(Q_0\)-module isometry onto \(p_0L^2(D_\mathcal U)\)
has initial projection \(q\in M_2(Q_0)\), of trace \(d'\).
Let its two basis entries be the row \(V\).
Initially these entries lie in \(L^2(D_\mathcal U)\), and
\begin{equation}
 Vq=V,\qquad E_{M_2(Q_0)}(V^*V)=q,\qquad
 \Theta(\gamma)V=Vb_\gamma ,
 \label{eq:scale-module-basis}
\end{equation}
where \(b_\gamma\) is a projective unitary representation
in \(qM_2(Q_0)q\).
The positive \(L^1\)-element \(VV^*=\sum_iV_iV_i^*\)
is invariant under \(\Theta\).  Its spectral projections
belong to the scalar relative commutant; hence
\begin{equation}
 VV^*=d'1.
 \label{eq:scale-bounded-basis}
\end{equation}
This also proves that every \(V_i\) is bounded.  Thus
\begin{equation}
 e_1=(d')^{-1}V^*V\le q
 \label{eq:scale-excess-projection}
\end{equation}
is a projection of trace one and commutes with all \(b_\gamma\).

\emph{Excluding excess module dimension.}
The projection $q$ has trace $d'$, whereas $e_1$ has trace one.
Thus $e_1=q$ will force $d'=1$ and turn the invariant module into
one unitary copy of $Q_0$. Suppose instead that $e_1\ne q$, and put
\[
 H_1=e_1-q/d'.
\]
Then $H_1\ne0$, $E_{M_2(Q_0)}H_1=0$, and $H_1$ commutes with
$b_\gamma$. This is the nonzero column whose long Fourier mass will
contradict the amplified gap.  Orthogonality to \(Q_0\) implies
that representing columns of \(H_1\) have vanishing Fourier
mass below every fixed multiple of \(L_n\).  Indeed test
orthogonality against the smooth positive radial cutoff
of the same bounded column; this cutoff belongs to \(Q_0\)
and is equal to one on the desired lower window.

Here and below the finite matrix trace only changes constants.
Choose bounded representatives \(G_n\) of \(H_1\) with exact
support \(r_T\ge l_n\), and short approximations to \(q,b_\gamma\)
with support \(r_T\le a_n\), so that
\begin{equation}
 L_n\ll a_n\ll d_n\ll l_n,\qquad
 \|G_n-H_{1,n}\|_2\longrightarrow0.
 \label{eq:scale-separated-scales}
\end{equation}
All approximations are uniformly operator-bounded.  Since these
are elements of fixed $2\times2$ matrix algebras over normalized
finite factors, they also satisfy
\[
 \|x_n\|_2\le\sqrt2\,\|x_n\|_\infty.
\]
Thus both norm bounds in Proposition~\ref{prop:ray-lift} and
Corollary~\ref{cor:ray-small-change} apply to every short or long
column used here.  One
explicit diagonal choice is as follows.  For the first \(j\)
group labels choose a fixed multiple \(M_jL_n\) that
approximates their tight columns within \(1/j\).  Require
the long column to have mass less than \(1/j\) below
\(j^2M_jL_n\), and pass to indices satisfying this and the
finitely many algebraic approximations.  Then use
\(a_n=M_jL_n\), \(d_n=jM_jL_n\),
\(l_n=j^2M_jL_n\), after harmless fixed slack in the cutoffs.
This construction produces an ordinary subsequence.

Lift the coisometric row \(V/\sqrt{d'}\) to exact
coisometric rows \(W_n\).  This is obtained by polar
decomposition, cutting a small spectral neighborhood of
zero, and filling the vanishing trace defects in the finite
factors.  The corrections tend to zero in Hilbert norm:
away from zero polar correction is uniformly continuous,
and the discarded support has trace tending to zero.
Put \(e_{1,n}=W_n^*W_n\) and
\[
 \Theta'_n(\gamma)=W_n^*\Theta_n(\gamma)W_n .
\]
These are exact projective representations in the
\(e_{1,n}\)-corners.  Choose projection lifts \(q_n\)
and corner-unitary lifts \(b_{\gamma,n}\).  The identities
in the ultraproduct give
\begin{equation}
 \|\Theta'_n(\gamma)-e_{1,n}b_{\gamma,n}\|_2\longrightarrow0.
 \label{eq:scale-corner-approximation}
\end{equation}
Neither the representatives \(b_{\gamma,n}\) nor their short
cuts need satisfy an exact representation law.

Let \(\xi_n=\mathbb V_{d_n}G_n\).
Proposition~\ref{prop:ray-lift}, first for short
approximations and then using the density error
\eqref{eq:ray-density-error}, gives
\begin{equation}
 \Delta_T(b_{\gamma,n})\xi_n-\xi_nb_{\gamma,n}\to0,\qquad
 \Delta_T(q_n)\xi_n-\xi_n\to0,\qquad
 \xi_nq_n-\xi_n\to0
 \label{eq:scale-lifted-intertwining}
\end{equation}
in Hilbert norm.  These statements use only the ordinary
identities \(b_\gamma H_1=H_1b_\gamma\) and \(qH_1=H_1q=H_1\).
Put \(\xi'_n=\Delta_T(e_{1,n})\xi_n\).
The left density of \(\xi_n\) has Fourier support in
\(\{r_T\le4d_n+C\}\), since the two ray tubes can overlap
only when their base points are at that distance.
Consequently
\[
 \langle\Delta_T(G_n)\xi_n,\xi_n\rangle=0
\]
for large \(n\).  The identity
\(e_{1,n}=G_n+q_n/d'+o_{L^2}(1)\) and the uniform density
bound now yield
\begin{equation}
 \|\xi'_n\|^2=(d')^{-1}\|\xi_n\|^2+o(1)
 \longrightarrow(d')^{-1}\|H_1\|_2^2>0 .
 \label{eq:scale-nonzero-lift}
\end{equation}
Using \eqref{eq:scale-corner-approximation} and
\eqref{eq:scale-lifted-intertwining}, we also obtain
\begin{equation}
 \Delta_T(\Theta'_n(\gamma))\xi'_n-\xi'_nb_{\gamma,n}\to0.
 \label{eq:scale-almost-intertwiner}
\end{equation}
Extend \(b_{\gamma,n}\) by the identity off \(q_n\).
This does not change the last assertion, since
\(\xi'_nq_n-\xi'_n\to0\).

The normalized left vector states defined by \(\xi'_n\)
on the semifinite corner with unit \(\Delta_T(e_{1,n})\)
are now almost invariant in predual norm.  Right
multiplication by the extended \(b_{\gamma,n}\) does not
change such a state; the error in
\eqref{eq:scale-almost-intertwiner} changes it by at most
twice the vector error times the vector norm.
Their positive density square roots are almost invariant
Hilbert unit vectors by the Powers--St\o rmer inequality
\cite{PowersStormer1970}, whose short trace proof we include:
\begin{equation}
 \|A^{1/2}-B^{1/2}\|_2^2\le\|A-B\|_1
 \quad(A,B\ge0).
 \label{eq:scale-square-root}
\end{equation}
To verify this inequality, put \(S=A^{1/2}\), \(T=B^{1/2}\),
and \(D=S-T\).  The trace duality inequality and cyclicity give
\[
 \|A-B\|_1\ge
 \Tr\bigl(\operatorname{sgn}(D)(S^2-T^2)\bigr)
 =\Tr\bigl(|D|(S+T)\bigr).
\]
Subtracting \(\Tr(D^2)\) from the last expression leaves
\(2\Tr(D_+T)+2\Tr(D_-S)\ge0\).
All products are trace-class by Hilbert--Schmidt
Cauchy--Schwarz; finite spectral truncation gives the same
calculation in the semifinite setting.
Conjugation on this corner is unitarily equivalent, via
\(\Delta_T(W_n)\), to the original amplified conjugation.
Admissibility and the fixed Kazhdan gap contradict these
almost invariant unit vectors.  This proves \(e_1=q\).

\emph{Recovering a near-identity unitary.}
Taking traces gives \(d'=1\). A matrix equivalence over
the factor \(Q_0\) reduces \(q\) to a single unit, and the
row \(V\) becomes an actual unitary \(u\) with
\(u^*\Theta(\Gamma)u\subset Q_0\) and
\(p_0=ue_{Q_0}u^*\).
Moreover
\begin{equation}
 \|E_{Q_0}u\|_2^2=\Tr(e_{Q_0}p_0).
 \label{eq:scale-jones-pairing}
\end{equation}
Align the polar part of \(E_{Q_0}u\) by a right unitary
from \(Q_0\); partial polar isometries extend to unitaries
in a finite factor.  Since \(E_{Q_0}u\) is a contraction,
the aligned unitary satisfies
\begin{equation}
 \|u-1\|_2^2
 \le2-2\|E_{Q_0}u\|_2^2
 =\|p_0-e_{Q_0}\|_2^2.
 \label{eq:scale-close-unitary}
\end{equation}
Lift \(u\) by unitaries and diagonalize the countably many
tightness conditions.  Choosing \(\varepsilon_K\) sufficiently
small gives the requested bound.  Inner conjugacy preserves
relative irreducibility; applying \(\Delta_T\) to the
conjugating unitary proves preservation of amplified
admissibility.
\end{proof}

\subsection{Composition of admissible kernels}

\begin{lemma}[Exact insertion]\label{lem:scale-insertion}
Let
\(\delta:L_\alpha(\Gamma)\to B_0\bar\otimes L_\beta(\Gamma)\)
be a trace-preserving embedding whose group representation is
relatively irreducible and has a sequence escaping in Fourier
probability in every ambient simple factor.  Let \(\Theta\)
be admissible for an isolated measured factor \(T\).
There is an admissible projective representation
\(\Theta\circledast\delta\), with source multiplier \(\alpha\),
whose spherical Fourier laws are the exact compositions of
the spherical Fourier kernels of \(\delta\) and \(\Theta\).
The final measured cocycle is the measured cocycle of \(\Theta\).
\end{lemma}

\begin{proof}
Write \(\xi\) for the source multiplier of \(\Theta\), and put
\[
 T_h(y)=\Theta(c(h,h^{-1}y)).
\]
The genuine section-cocycle identity gives
\begin{align*}
 T_g(y)T_h(g^{-1}y)
 &=\nu(g,h)(y)T_{gh}(y),\\
 \nu(g,h)(y)
 &=\xi\bigl(c(g,g^{-1}y),
           c(h,h^{-1}g^{-1}y)\bigr).
\end{align*}
Thus \(\nu\) is a normalized scalar-function-valued cocycle
for the action on \(Y\).  In the finite twisted crossed
product for this action use shift unitaries \(b_h\) with
multiplier \(\beta/\nu\).  Tensor this crossed product with
the final finite algebra \(D\).  Then
\begin{equation}\label{eq:scale-insertion}
 \Psi(w_h)=T_hb_h,\qquad
 \Psi(w_g)\Psi(w_h)=\beta(g,h)\Psi(w_{gh}).
\end{equation}
Orthogonality of the shifts makes \(\Psi\) trace-preserving
on the twisted group algebra.  It extends normally:
trace preservation gives equality of every even moment
of the absolute value of a group polynomial, hence equality
of its operator norm; the resulting trace-preserving
isometric map extends by bounded strong approximation.
Apply \(\id_{B_0}\otimes\Psi\) to \(\delta\).

Here is the kernel formula.  If
\(\delta(w_\gamma)=\sum_h a_{\gamma,h}\otimes w_h\), then,
for a Borel set \(E\) of final Fourier lengths, its mass
after insertion is
\begin{equation}\label{eq:scale-kernel-composition}
 \sum_{h\in\Gamma}\|a_{\gamma,h}\|_2^2
       \int_Y\bigl\|D_{1_{\{r_T\in E\}}}
                    \Theta(c(h,h^{-1}y))\bigr\|_2^2\,dy .
\end{equation}
Distinct \(h\)'s are orthogonal in the middle shift
coordinate.  Integration over a source spherical
transporter gives the same formula for spherical laws.
The change \(y\mapsto h^{-1}y\) identifies the inner
integral with the spherical kernel at the translation
\(h\in G\).  In particular all inserted kernels are
probability kernels.

It remains to prove admissibility and factoriality of
the coefficient tensor, rather than assume them.
Use \(\Psi(w_h)\) as a right Fourier basis for the middle
crossed product.  Its coefficient Hilbert space is
\[
 \mathcal E=L^2(Y;L^2(D)).
\]
Left multiplication by \(\Psi(w_g)\) acts on this coefficient
space by
\begin{equation}
 (\mathcal U_gk)(y)=\Ad T_g(y)\,k(g^{-1}y),
 \label{eq:scale-induced-coefficients}
\end{equation}
and shifts the basis from \(h\) to \(gh\), with multiplier
\(\beta(g,h)\).  Right multiplication only shifts the
basis from \(h\) to \(hg\), with multiplier
\(\beta(h,g)\).  The coefficient action is the restriction
of the \(G\)-induced trace action of \(\Ad\Theta\).
Its invariant space is \(\mathbb C1\).  On
\(\mathcal E\ominus\mathbb C1\), coefficients tend to zero
along sequences escaping in every simple factor, by
Howe--Moore and the product invariant-space decomposition
of Section~\ref{sec:setup}.

To apply this decay to Fourier-distributed unitaries, take
bounded finite basis tests \(b\,k\,\Psi(w_h)\) and
\(b'\,k'\,\Psi(w_{h'})\), and write
\(A=\sum_i a_iw_i\) for one of the escaping unitaries of
\(\delta\).  The conjugation pairing is bounded, up to
the fixed norms of \(b,b'\), by
\begin{equation}\label{eq:scale-composition-decay}
 \sum_{i,j:\,ih=h'j}
 \|a_i\|_2\|a_j\|_2
       |\langle k',\mathcal U_i k\rangle|.
\end{equation}
The relation determines \(j\) uniquely from \(i\).
Its sum without the last factor is at most
\(\sum_i\|a_i\|_2^2=1\), by Cauchy--Schwarz.
Outside fixed large compact windows in every factor,
the last coefficient is uniformly small.  The remaining
terms tend to zero by Fourier escape and Cauchy--Schwarz;
fixed shifts in the relation do not change escape.
Thus the conjugation pairing tends to zero on
the nonscalar coefficient summand.  On the scalar
coefficient summand it is precisely the conjugation
representation of \(\delta\), whose invariant space is
scalar.  Density proves scalar relative commutant for
the inserted representation.  The center of its ambient
finite tensor lies in that commutant, so the ambient
tensor is a factor.

For the amplified assertion replace \(\mathcal E\) by
\[
 \mathcal E_{\rm amp}=
 L^2\bigl(Y;
   L^2(B(L^2(X_T))\bar\otimes D,\Tr\otimes\tau)\bigr)
\]
and replace \(T_g\) by \(\Delta_T(T_g)\).
This induced coefficient representation has no invariant
vector, by admissibility of \(\Theta\).  Finite-rank
trace-Hilbert coefficient tests and the calculation
\eqref{eq:scale-composition-decay} now apply on the
entire amplified space.  They prove absence of invariant
vectors there.  The coefficient algebra of the resulting
tensor consists of \(B_0\), the middle finite crossed
product, and the coefficient factor of \(D\); its final
group factor is still the original measured factor of
\(\Theta\).  This proves all assertions.
\end{proof}

\subsection{A linear subsequence}

\begin{lemma}\label{lem:scale-many-increments}
Let \(\Theta_n\) satisfy the assumptions of
Proposition~\ref{prop:ray-addition}.  Suppose \(t_n\to\infty\)
and for some \(a,\alpha>0\),
\[
 P_{\Theta_n,R_n}(t_n)([a,\infty])\ge\alpha.
\]
If \(u_n/t_n\to\infty\), then
\begin{equation}
 P_{\Theta_n,R_n}(u_n)\longrightarrow\delta_\infty .
 \label{eq:scale-many-increments}
\end{equation}
This assertion is uniform under compact pre- and
post-translations.
\end{lemma}

\begin{proof}
Fix an integer \(m\).  For large \(n\),
\(u_n=mt_n+(u_n-mt_n)\) with nonnegative remainder.
Pass to a subsequence on which the laws of these
segments converge.  Repeated use of
Proposition~\ref{prop:ray-addition} makes the limit at
\(u_n\) equal to \(P^{*m}*Q\), where
\(P([a,\infty])\ge\alpha\) and \(Q\) is supported in
\([0,\infty]\).  For a finite \(M\),
\begin{equation}
 (P^{*m}*Q)([0,M])
 \le \mathbb P\{\operatorname{Bin}(m,\alpha)\le M/a\}.
 \label{eq:scale-binomial}
\end{equation}
The right side tends to zero as \(m\to\infty\).
Diagonal extraction over the countably many \(m\)'s
justifies taking them after the subsequential limit.
This proves \eqref{eq:scale-many-increments}.
Compact changes are covered by
Lemma~\ref{lem:ray-compactshort}; a failure of uniformity
would give a violating deterministic sequence.
\end{proof}

\begin{proposition}[Positive mass at linear speed]\label{prop:scale-linear-subsequence}
For each isolated source type \(i\), with matched measured
type \(p(i)\), there exist \(C<\infty\), \(\alpha>0\), and
\(T_n\to\infty\) such that, in the original regular module
model, every fixed \(\xi\in\mathcal H\) satisfies
\begin{equation}\label{eq:scale-linear-subsequence}
 \liminf_{n\to\infty}\ \inf_{T_n/2\le t\le T_n}
 \bigl\|D_{1_{\{r_{p(i)}\le CT_n\}}}U_{a_t}\xi\bigr\|^2
 \ge\alpha\|\xi\|^2 .
\end{equation}
The assertion also holds in normal Hilbert amplifications
and after compact source pre- and post-translations.
It makes no assertion uniform over a moving Hilbert unit
ball.
\end{proposition}

\begin{proof}
Use the representation \(\delta\) of
Lemma~\ref{lem:scale-original-model}.  Suppose first
that, on a pure isolated source type \(i\), its matched
length is totally superlinear:
\begin{equation}\label{eq:scale-superlinear-assumption}
 P_{\delta,t}(t)([0,C])\longrightarrow0
 \qquad(C<\infty).
\end{equation}
The contradiction will use the entire permutation cycle through $i$.
Composing the original representation around
\[
 i\longmapsto p(i)\longmapsto\cdots\longmapsto i
\]
returns to the same source and measured type. Under the assumption
\eqref{eq:scale-superlinear-assumption}, its first stage increases
length faster than linearly; the entropy bound prevents the later
stages from losing that increase. We will insert this cycle before
an admissible representation that reaches a large Fourier threshold
in nearly the least possible source time. The composition would
reach the same relative threshold in a strictly shorter time.

\emph{Choosing a nearly earliest witness.}
Fix numbers \(D'_n\to\infty\). For each \(n\), minimize
over all admissible representations with both source
and measured type equal to \(i\), all cutoffs \(L\ge1\),
and times \(t\ge1\), subject to
\begin{align}
 \|D_{1_{\{r_i>L\}}}\Theta(\gamma)\|_2
 &\le\varepsilon_K
       &&(\gamma\in\mathcal F),\label{eq:scale-minimal-first}\\
 P_{\Theta,D'_nL}(t)([1,\infty])
 &\ge\tfrac12 .&&\label{eq:scale-minimal-second}
\end{align}
Here $L$ controls the Fourier tails of the fixed Kazhdan set, and
$D'_nL$ is the output threshold to be reached. Choose a witness
\(t_n\) at most twice the infimum, and write its cutoff as
\(L_n\) and its output scale as \(R_n=D'_nL_n\).

The feasible set is nonempty.  Use the identity group
embedding into \(L_\eta(\Gamma)\), with any fixed cocycle,
and take \(L\) containing the finite Kazhdan set in the
measured length.  Relative irreducibility is ICC;
amplified admissibility follows from all-factor escape
of lattice elements and translation coefficient decay.
For this embedding the transporter length on a pure
axis differs from the source length by a bounded amount
outside arbitrarily small section sets.  Thus sufficiently
large \(t\) meets \eqref{eq:scale-minimal-second}.

\emph{Tightening the witness and repeating increments.}
Apply Lemma~\ref{lem:scale-tightening}, choosing
\(\varepsilon_K\) so small that the resulting conjugacies
preserve at least \(1/4\) mass in
\eqref{eq:scale-minimal-second}.  This preservation is
uniform on moving transporters: for every unitary \(x\),
\(\|u_n^*xu_n-x\|_2\le2\|u_n-1\|_2\), and the same
bound holds after integration over the section cocycle.
Rename the tightened representations \(\Theta_n\).
Each fixed element is tight at scale \(L_n\), hence
short on scale \(R_n\).

After extraction \(t_n\to\infty\).  Otherwise
Lemma~\ref{lem:ray-compactshort}, with its proof applied to
tightness on scale \(L_n\), would make all bounded-time
transporters short on scale \(R_n=D'_nL_n\), contradicting
the retained \(1/4\) mass.  Lemma~\ref{lem:scale-many-increments}
therefore gives
\begin{equation}\label{eq:scale-minimal-escape}
 u_n/t_n\to\infty
 \quad\Longrightarrow\quad
 P_{\Theta_n,R_n}(u_n)\longrightarrow\delta_\infty .
\end{equation}

\emph{Shortening the witness by a cycle of insertions.}
Now insert copies of \(\delta\) before \(\Theta_n\), once
for each step around the finite permutation cycle
through \(i\), using Lemma~\ref{lem:scale-insertion}.
The first insertion, viewed from the source, is the
totally superlinear one in
\eqref{eq:scale-superlinear-assumption}.  Choose
\(m_n\to\infty\) sufficiently slowly that
\begin{equation}
 t'_n=t_n/m_n\to\infty,\qquad
 P_{\delta,t_n}(t'_n)([m_n,\infty])\longrightarrow1 .
 \label{eq:scale-small-start}
\end{equation}
Such a choice follows directly from total superlinearity:
for each integer \(m\), choose a threshold after which
the ratio of measured length to source length exceeds
\(m^2\) with probability at least \(1-1/m\), and require
\(t_n/m\) to exceed that threshold.

Every subsequent insertion retains length much larger
than \(t_n\), by Lemma~\ref{lem:scale-entropy}.  Its
complementary coordinates remain tight.  To justify
this assertion for the random middle indices in
\eqref{eq:scale-kernel-composition}, first restrict
the complementary coordinates to a fixed compact
window with arbitrarily large probability.  Cartan
decomposition then writes those indices as pure
translations with compact pre- and post-factors.
The entropy estimate is uniform on these windows.
After the finitely many insertions, the terminal
kernel \(\Theta_n\) is tested on its source type \(i\)
at lengths divided by \(t_n\) tending to infinity in
probability.  Equation~\eqref{eq:scale-minimal-escape}
and its compact uniformity show that the final measured
length divided by \(R_n\) tends to infinity in probability.

For a fixed source group element the preceding finite
chain of kernels instead gives a fixed probability
measure on the countable set of terminal input labels.
Its mass can be restricted to a finite set with
arbitrarily small loss.  Tightness of the terminal
fixed-label laws at scale \(L_n\), and the uniform
section approximation for each such label, prove
tightness at scale \(L_n\) for every fixed group
element of the composite representation.  Hence,
for a fixed sufficiently large \(C_1\), and after
extraction, the composite satisfies
\eqref{eq:scale-minimal-first} with cutoff \(C_1L_n\).
Its final escape also gives
\eqref{eq:scale-minimal-second} with threshold
\(C_1R_n\), at time \(t'_n\).
The ratio of these new scales is still \(D'_n\).
Admissibility is preserved by every insertion.
Thus this is a feasible witness in the very same
minimization, and
\[
 t_n\le2\inf t\le2t'_n,
\]
contrary to \(t'_n/t_n\to0\).

\emph{From a scalar subsequence to the fixed-vector conclusion.}
Total superlinearity is therefore impossible. There
are \(C_0,\alpha_0>0\) and \(T_n\to\infty\) for which
the scalar law has mass at least \(\alpha_0\) below
\(C_0T_n\).  This is uniform for
\(t\in[T_n/2,T_n]\), after enlarging \(C_0\) slightly.
Indeed, a violating sequence \(t_n\) would, by
Proposition~\ref{prop:ray-addition} applied to the fixed
representation \(\delta\) on scale \(T_n\), give a
convolution limit for \(T_n=t_n+(T_n-t_n)\).
The mass of a nonnegative sum below \(C_0\) cannot
exceed the mass of either summand below \(C_0\).
Weak-limit boundary issues are removed by the stated
slack in the cutoff.

Lemma~\ref{lem:links-invariant-variables} identifies all
these limiting effects as scalar.  The corner transfer
in Lemma~\ref{lem:scale-original-model} identifies
their scalar values in the original regular model.
It follows that the same lower bound holds on every
fixed bounded column, uniformly over the indicated
times.  Effect boundedness and Hilbert density extend
it to every fixed \(\xi\), proving
\eqref{eq:scale-linear-subsequence}.  Tensor density
gives normal Hilbert amplifications.  The compact
perturbation statement follows from the same short
approximation and rescaled-effect argument.
\end{proof}

\subsection{Complementary localization and the opposite-pair estimate}

Fix an isolated source factor $S$ and its matched target factor $T$.
The two factors occur in copies of the same ambient group $G$; the
arbitrary right group $\Lambda$ supplies no geometric factor. Write
$\widehat G_S$ and $\widehat G_T$ for the respective complementary
products. We work first with one measured copy. Its pure-vertex
family is denoted by $x\mapsto P_x$, for $x\in B_S=\partial S$,
with values in the projection-valued measures on $B_T=\partial T$.
The sharpness and covariance from Section~\ref{sec:transfer} give
\[
 \sigma_sP_x=P_{sx},\qquad
 \beta_hP_x(f)=P_x(f\circ h^{-1}),
\]
and the complementary source factors fix this family. Let $d_S,d_T$
be the compact Carnot distances, let $dx$ be compact angular probability
on $B_S$, and let $Q_S,Q_T$ be the homogeneous dimensions.

Our next goal is to control how the measurements change when their source
labels approach one another. We will integrate the squared norms of
ordered products $P_y(g)P_x(f)$, where $f,g$ are smooth target functions
with separated supports.
Lattice packing gives a bound proportional to $1+L$ when measured Fourier
length is at most $L$. The stabilizer of an opposite source pair has a
split parameter. For suitable $T_n\to\infty$ and $A_0<\infty$, its
positive segment of length comparable to $T_n$ retains positive Fourier
mass below $L=A_0T_n$ for each fixed input vector, by
Proposition~\ref{prop:scale-linear-subsequence}. Averaging over that segment
and dividing by $T_n$ removes the length factor from the packing bound.
Unfolding then gives the integrated opposite-pair estimate used for
horizontal differences in the next subsection.

Every lattice count in this argument keeps a compact window in
$\widehat G_T$. The projection of $\Gamma$ to $T$ need not be discrete.
We first construct an isometry that permits this complementary window
while preserving the source action and the angular measurements.

\begin{lemma}[Complementary localization]
\label{lem:scale-complement-isometry}
There is a base-decomposable isometry
\[
 V:\mathcal H\longrightarrow
 \ell^2\otimes L^2(\widehat G_T)\otimes\mathcal H
\]
such that, for $s\in S$ and measured labels $(h,\lambda)$,
\begin{align*}
 VU_s&=(1\otimes1\otimes U_s)V,\\
 VP_x(f)&=(1\otimes1\otimes P_x(f))V,\\
 VW_{(h,\lambda)}
  &=(1\otimes L_{\hat h}\otimes W_{(h,\lambda)})V.
\end{align*}
Here $L_q\zeta(p)=\zeta(q^{-1}p)$, and $\hat h$ is the complementary
projection of $h$.  Consequently, for every compact
$C\subset\widehat G_T$, the positive contraction
\[
 W_C=V^*(1\otimes1_C\otimes1)V
\]
commutes with the base multipliers, $U_S$, $W_\Lambda$, and all the
$P_x(f)$.  For a compact exhaustion, $W_C\to1$ strongly.
\end{lemma}

\begin{proof}
Use the regular corner from the complementary regular realization in Section~\ref{sec:setup} and
amplify its left Fourier action by
$\widehat\Delta(w_{(h,\lambda)})=L_{\hat h}\otimes w_{(h,\lambda)}$.
In the rectangular trace Hilbert space with corners
$\widehat\Delta(e),e$, source transport acts by
\[
 Z(y)\longmapsto
 \widehat\Delta(X_s(s^{-1}y))Z(s^{-1}y)X_s(s^{-1}y)^*.
\]
The scalar projective phases cancel.  We show that every nonzero
projection $d_2\in\mathcal N^S$ contains the right support of a nonzero
invariant column.

Choose a unit vector $\psi\in L^2(\widehat G_T)$ with strictly positive
translation coefficient
$q\mapsto\langle L_q\psi,\psi\rangle$, bounded below on each compact
set.  For example, choose an everywhere positive square-integrable
function; its coefficient is positive and continuous.  Test the
preceding action on
$\widehat\Delta(d_2)(\psi\otimes d_2)$.  Reindexing the base and using
$X_s(y)d_2(y)X_s(y)^*=d_2(sy)$ expresses its overlap with its $s$-translate
as the sum of squared Fourier coefficients of $X_s(y)d_2(y)$ weighted
by $\langle L_{\hat h}\psi,\psi\rangle$.  The total unweighted mass is
$\|d_2\|_2^2$.  Complementary tightness for a pure isolated type makes
these overlaps bounded below by a positive constant, uniformly in
$s\in S$.  The uniformity follows first on a Cartan ray and then on
all of $S$ by Cartan decomposition and the compact-transport
approximation used above; bounded
Cartan parameters form a strongly continuous compact family.

The element of least norm in the closed convex hull of this orbit is
therefore nonzero and invariant.  Its right support is contained in
$d_2$.  Polar decomposition in the rectangular affiliated-operator
space gives an invariant partial-isometry column; its initial
projection belongs to $\mathcal N^S$.  Repeat in the orthogonal
complement of the initial projection.  Separability gives a countable
exhaustion.  Place the resulting columns in separate multiplicity
coordinates.  Their direct sum $V_\Delta$ satisfies
$V_\Delta^*V_\Delta=e$ and the source intertwining relation.  In
particular it is bounded, and its integrated rectangular trace
Hilbert norm squared is $(\Tr\otimes\tau)(e)=d$.

We now remove the left amplification.  For a Fourier output index
$k$ and parameter $P\in\widehat G_T$, use the unitary change of variables
\[
 F_k(P)\longmapsto \widetilde F_k(p)=F_k(\hat k p).
\]
Left multiplication by $\widehat\Delta(w_h)$ becomes ordinary left
multiplication, while right multiplication by $w_h^*$ acquires
$L_{\hat h}$ in the new variable.  This proves the source and right
action relations for the resulting isometry $V$.

The same change of variables leaves diagonal angular Fourier tests
unchanged.  To pass such a test through $V_\Delta$ on a defining pure
path, first trim the extra Hilbert coordinate to finite rank and
approximate the remaining column by bounded finite Fourier columns
in integrated trace Hilbert norm.  A fixed left Fourier shift does
not change the limiting image direction.  Moving each such shift
through the angular test leaves its uncut Gram factor.  The
approximation error tends to zero on bounded input columns because
$V_\Delta^*V_\Delta=e$; Hilbert density then handles every input.
Thus compression of the constant-amplified limiting angular measure
by $V$ is $P_x$.  Both measures are sharp.  For a spectral projection
$E$ with $V^*EV=P$, the equality
$\|(EV-VP)\xi\|^2=0$ proves exact intertwining.  This proves the
angular relation.  The assertions about $W_C$ now follow directly
from the displayed intertwinings and monotone exhaustion of
$L^2(\widehat G_T)$.
\end{proof}

\begin{definition}[Controlled angular masks]
\label{def:scale-mask-controls}
Fix an integer $m\ge1$ and positive constants specifying support
diameter, separation, and overlap.  A controlled family at scale
$0<\delta\le\delta_0$ consists of smooth functions
$\{f_A\}_A,\{g_D\}_D$ on $B_T$ and a set $\mathscr R$ of ordered pairs
$(A,D)$.  Each support lies in a ball of radius $a\delta$; each family
of centers has boundedly many points in every ball of radius
$b\delta$, for each fixed $b$; and
\[
 \operatorname{dist}_T(\supp f_A,\supp g_D)\ge c\delta
 \quad ((A,D)\in\mathscr R).
\]
After translation and graded dilation to a unit chart, derivatives
through order $m$ are bounded by one.  Changing this last bound only
changes the constants by scalar rescaling.  We permit any fixed
enlargement of the diameter, overlap, or separation parameters.
Set
\[
 p_\delta(x,y)
 =\sum_{(A,D)\in\mathscr R}
   P_x(f_A)^*P_y(|g_D|^2)P_x(f_A).
\]
This is a positive operator.  Bounded overlap and spectral calculus
give $0\le p_\delta(x,y)\le C_0 1$, with $C_0$ depending only on the
controls, without any commutation assumption between $P_x$ and $P_y$.
\end{definition}

\begin{lemma}[Packing with a complementary window]
\label{lem:scale-pair-packing}
There is an integer $m_0$, depending only on the fixed boundary
charts and the geometric support, separation, and overlap parameters, such that for each compact
$C\subset\widehat G_T$, controlled masks of order $m\ge m_0$ satisfy
\begin{equation}
 \sum_{\substack{h\in\Gamma:\ r_T(h)\le L,\ \hat h\in PC^{-1}}}
 W_h p_\delta(x,y)W_h^*
 \le C_C(1+L)\delta^{-Q_T}1
 \label{eq:input-9}
\end{equation}
for every $P\in\widehat G_T$, $L\ge1$, and source opposite pair
$(x,y)$.  The constant is independent of $P,L,\delta,x,y$ and the
number of masks.
\end{lemma}

\begin{proof}
We prove a bound $C_C\delta^{-Q_T}$ in each unit length shell.  Choose
a sufficiently small relatively compact identity neighborhood $U$
such that $hU$, $h\in\Gamma$, are disjoint.  Right thickening changes
$r_T(h)$ by a bounded amount and changes the forward direction $h^+$
by $O(e^{-r})$ in a shell $r\le r_T(h)<r+1$.  It enlarges
$PC^{-1}$ to a translate of a fixed compact set.  Cartan integration
therefore gives, for a ball of radius $R\ge e^{-r}$,
\begin{equation}
 \#\{h:r\le r_T(h)<r+1,\ \hat h\in PC^{-1},\ h^+\in B(z,R)\}
 \le C_C(1+Re^r)^{Q_T}.
 \label{eq:scale-shell-count}
\end{equation}
Indeed the angular ball has measure at most
$C(R+e^{-r})^{Q_T}$, the radial shell has Jacobian at most $Ce^{Q_Tr}$,
and the complementary factor has bounded Haar volume independent of
its left translate $P$.  This argument counts full lattice points,
not points in a possibly nondiscrete projection.

We record the two boundary estimates used below.  Put
$R_0=C'e^{-r}/\delta$.  A set at distance at least $c\delta/2$ from
$h^-$ has image under $h$ contained in $B(h^+,R_0)$.  Moreover, if
$\supp f_A$ is within $c\delta/2$ of $h^-$, then on an annulus at
distance $\rho\ge R_0$ from $h^+$,
\begin{equation}
 \|f_A\circ h^{-1}-f_A(h^-)\|_{C^m(\rho\text{-normalized charts})}
 \le C_m\frac{R_0}{\rho}.
 \label{eq:scale-pole-remainder}
\end{equation}
Here one may enlarge $C'$ by a fixed factor.  Use the normalized
Carnot charts of Lemma~\ref{lem:ray-geometry}.  To see these estimates,
put the two polar endpoints at zero and infinity.  The Cartan element
is a graded dilation of ratio $e^{-r}$ in the chart at the attracting
endpoint.  Chart inversion sends an annulus of radius $\rho$ to an
annulus of radius comparable to $\rho^{-1}$, with uniformly bounded
derivatives after graded normalization.  Thus the inverse image of
the annulus is within $O(e^{-r}/\rho)$ of the repelling endpoint.
Taylor's formula for a mask whose derivatives are normalized at
scale $\delta$ proves \eqref{eq:scale-pole-remainder}.  Every
nonconstant normalized derivative has the same bound: its chain-rule
factors contain at least one positive-degree dilation factor.
Compact polar rotations preserve the chosen horizontal metric, so
all constants are uniform in $h$.

If $R_0$ is bounded below by a fixed positive number, total shell
counting already gives $C_C\delta^{-Q_T}$.  Assume otherwise.  From
\eqref{eq:scale-shell-count}, for either one of the individual PVMs
$P_x,P_y$, we obtain
\begin{equation}
 \sum_h P_z(1_{B(h^+,bR_0)})\le C_{C,b}\delta^{-Q_T}1
 \quad (z=x\text{ or }y),
 \label{eq:scale-pole-projection-sum}
\end{equation}
where the sum is over the shell and complementary window.  This is
just the pointwise counting inequality integrated against a single
PVM.

Conjugation by $W_h$ composes the masks with $h^{-1}$.  For those
$A$ whose supports are at distance at least $c\delta/2$ from $h^-$,
the rightmost masks are supported in $B(h^+,R_0)$.  Bounded overlap
gives
\begin{align*}
 &\sum_{A,D}
 \|P_y(g_D\circ h^{-1})P_x(f_A\circ h^{-1})\xi\|^2\\
 &\hspace{20mm}\le
 C\sum_A\|P_x(f_A\circ h^{-1})\xi\|^2
 \le C\|P_x(1_{B(h^+,R_0)})\xi\|^2.
\end{align*}
Summing in $h$ uses \eqref{eq:scale-pole-projection-sum} for $P_x$.

Only boundedly many remaining $A$ occur for each $h$.  Separation
implies that all their paired $g_D$ are far from $h^-$.  Summing the
$g_D$ squares reduces their contribution to
\[
 C\sum_{h,A\text{ remaining}}
 \|P_y(1_{B(h^+,R_0)})P_x(f_A\circ h^{-1})\xi\|^2.
\]
Subtract the scalar $f_A(h^-)$.  The scalar term is controlled by
\eqref{eq:scale-pole-projection-sum} for $P_y$.  The remainder in an
enlarged pole ball is controlled by that equation for $P_x$, using
its uniformly bounded amplitude.  Decompose the rest into annuli of
radii $\rho_j=2^jR_0$.  On each annulus,
\eqref{eq:scale-pole-remainder} gives normalized size $O(2^{-j})$.

Here is the summation detail that preserves the order of the two
PVMs.  Fix a bounded-overlap global cover by cells of scale $\rho_j$.
Only boundedly many cells meet any one of these annuli.  Smooth
expansion on each normalized cell gives
\[
 F_{h,A,j}=
 \sum_{E,\nu}c_{h,A,E,\nu}\varphi_{E,\nu},\qquad
 |c_{h,A,E,\nu}|
 \le C_m2^{-j}(1+|\nu|)^{-m},
\]
where $F_{h,A,j}$ is the annular remainder,
$\varphi_{E,\nu}$ is supported in an enlargement of $E$, and the
functions $\varphi_{E,\nu}$ do not depend on $h$ or $A$.  For fixed
$\nu$ they have a uniformly bounded square sum over $E$.  The
expansion follows by a smooth cutoff to a unit cube and its Fourier
series; taking $m$ above the chart dimension makes the weights
summable.  Cauchy--Schwarz over the bounded number of cells for a
fixed $h,A$ and then \eqref{eq:scale-pole-projection-sum} for $P_y$
give
\begin{align*}
 &\sum_{h,A,E}|c_{h,A,E,\nu}|^2
 \|P_y(1_{B(h^+,R_0)})P_x(\varphi_{E,\nu})\xi\|^2\\
 &\qquad\le
 C_m2^{-2j}(1+|\nu|)^{-2m}\delta^{-Q_T}
 \sum_E\|P_x(\varphi_{E,\nu})\xi\|^2\\
 &\qquad\le
 C_m2^{-2j}(1+|\nu|)^{-2m}\delta^{-Q_T}\|\xi\|^2.
\end{align*}
In the first inequality the vector $P_x(\varphi_{E,\nu})\xi$ is
independent of $h$.  Sum first over frequencies and then over
annuli in the Hilbert norm of the $(h,A)$-indexed family.  The sums
$\sum_\nu(1+|\nu|)^{-m}$ and $\sum_j2^{-j}$ converge.  This proves
the required shell bound.  Summing the $O(1+L)$ shells proves
\eqref{eq:input-9}.
\end{proof}

Let $B_S^{(2)}$ be the space of ordered opposite pairs.  In rank one
this is $B_S\times B_S$ with the diagonal removed.  Let $dm$ be its
$S$-invariant measure.  In compact charts,
\begin{equation}
 dm(x,y)\asymp d_S(x,y)^{-2Q_S}\,dx\,dy.
 \label{eq:scale-pair-measure}
\end{equation}
For completeness, first use horospherical Haar measure $dn$ in a
source boundary chart. Write the smooth positive pair density there as
$F(n,n')\,dn\,dn'$. Simultaneous horospherical translation gives
$F(n,n')=F(e,n^{-1}n')$. A split dilation scales each Haar variable
by $t^{Q_S}$ and preserves the pair measure, so
\[
 F(e,\delta_t w)=t^{-2Q_S}F(e,w).
\]
The density is smooth and strictly positive on a compact homogeneous
unit annulus. It is therefore comparable to the homogeneous gauge to
the power $-2Q_S$. On buffered compact charts, compact angular
probability has a smooth density bounded above and below relative to
$dn$, and the Carnot distance is comparable to that gauge. Finitely
many such charts give \eqref{eq:scale-pair-measure}.

\begin{proposition}[Integrated opposite-pair estimate]
\label{prop:scale-pair-estimate}
There is a dense linear subspace $\mathscr E\subset\mathcal H$ such
that, for each fixed set of mask controls and each $\xi\in\mathscr E$,
there is a finite constant $C_\xi$ with
\begin{equation}
 \int_{B_S^{(2)}}
  \langle\xi,p_\delta(x,y)\xi\rangle\,dm(x,y)
 \le C_\xi\delta^{-Q_T}
 \label{eq:input-10}
\end{equation}
for all controlled families and $0<\delta\le\delta_0$.
The subspace can be chosen invariant under the source and target
actions and under every bounded operator commuting with all
$P_x(f)$.  Constants may change with the vector and the fixed
controls, but not with $\delta$ or the chosen masks.
\end{proposition}

\begin{proof}
Fix an opposite pair $(x_0,y_0)$ and let $J_0=MA<S$ be its
compact-by-split stabilizer. Temporarily write
$p_\delta=p_\delta(x_0,y_0)$. We first place a Parseval frame below this
positive effect, then apply Fourier packing to its complementary
localization. The stabilizer average will cancel the length factor,
after which we can unfold the estimate over all opposite pairs.

Choose complementary Parseval frame columns $b_j\in H$ for the joint
action $\pi_0,W_\Lambda$; thus
\[
 \sum_{\gamma,\lambda,j}
 |W_\lambda\pi_0(\gamma)b_j\rangle
 \langle W_\lambda\pi_0(\gamma)b_j|=1
\]
in the strong sense. Such a finite frame of orbit generators is part of
the complementary regular realization. Let $f\ge0$ be bounded and
compactly supported on $G$, with
\begin{equation}
 \int_{J_0}f(ag)\,da\le1\qquad(g\in G).
 \label{eq:scale-fibre-density}
\end{equation}
Using the induction section $d_0$, define the fixed Hilbert vectors
\[
 v_{\gamma j}(o')=
 f(d_0(o')\gamma)^{1/2}
 p_\delta(o')^{1/2}\pi_0(\gamma)b_j.
\]
The vectors depend on the masks and on $f$, which remain fixed throughout
the later limit in $n$. Since $J_0$ fixes the reference pair, $U_a$
commutes with $p_\delta$. At each base point, reindexing $\gamma$ by the
section cocycle and using \eqref{eq:scale-fibre-density} gives the
following inequality of positive operators on the fibre $H$:
\[
 \int_{J_0}\sum_{\lambda,\gamma,j}
 |W_\lambda(U_av_{\gamma j})(o')\rangle
 \langle W_\lambda(U_av_{\gamma j})(o')|\,da
 \le p_\delta(o')
 \quad\text{for almost every }o'.
\]
Its global rank-one consequence on $\mathcal H=L^2(Y;H)$ is
\begin{equation}
 \int_{J_0}\sum_{\lambda,\gamma,j}
 |W_\lambda U_av_{\gamma j}\rangle
 \langle W_\lambda U_av_{\gamma j}|\,da
 \le p_\delta.
 \label{eq:scale-stabilizer-frame}
\end{equation}
Indeed the reindexed scalar is $f(a^{-1}d_0(o')\gamma')$, where
$\gamma'=c(a,a^{-1}o')\gamma$; its integral is at most one. The
remaining sum is the Parseval identity compressed by
$p_\delta^{1/2}$. We retain this pointwise positive inequality for the
Fourier-row tests below. It also gives the displayed inequality
on the global Hilbert space by first applying
\[
 \left|\int_Y\langle\xi(y),v(y)\rangle\,dy\right|^2
 \le \int_Y|\langle\xi(y),v(y)\rangle|^2\,dy
\]
and then summing and integrating. The projective phases disappear in
these pointwise squared coefficients.

We next identify the Fourier norm to which packing applies. Let
$V_\Delta$ be the column in the proof of
Lemma~\ref{lem:scale-complement-isometry}, before the change of variables
that removes its left amplification. At complement parameter $P$ and
base point $o'$, an adjoint identity row is
$V_\Delta(P,o')^*\mathbf1_{\rm row}$. For a fibre input $z\in H$,
its pairing with $W_hW_\lambda z$ reads the corresponding output row of the
$(h,\lambda)$ Fourier coefficient of $V_\Delta z$:
\[
 \left\langle V_\Delta(P,o')^*\mathbf1_{\rm row},
                 W_hW_\lambda z\right\rangle.
\]
This uses the right-adjoint convention for $W$. Summing the squared
coefficients over rows and $\lambda$ is therefore the Fourier Hilbert
norm in the measured $h$ coordinate. After the change of variables in
the localization lemma, $P=\hat h p$. Thus
\[
 p\in C\quad\Longleftrightarrow\quad \hat h\in PC^{-1}.
\]
This is precisely the complementary window in the packing inequality.

Test \eqref{eq:input-9} on these adjoint rows, insert the pointwise
frame inequality underlying \eqref{eq:scale-stabilizer-frame}, and
integrate in $P,o'$. The integrated
sum of their squared trace norms is
$(\operatorname{Tr}\otimes\tau)(e)=d$. The preceding coefficient identity
and Parseval therefore give, with $D_{\le L}$ the measured Fourier
projection onto $r_T(h)\le L$,
\begin{equation}
 C_C(1+L)\delta^{-Q_T}
 \ge\int_{J_0}\sum_{\gamma,j}
       \|D_{\le L}1_CVU_av_{\gamma j}\|^2\,da.
 \label{eq:scale-stabilizer-length-bound}
\end{equation}
The coefficient calculation is first made for bounded columns of finite
Fourier support. Approximation and the finite integrated squared norm
of the rows extend it to the fixed Hilbert vectors above. Positivity
permits all sums and integrals by Tonelli's theorem.

We now choose the stabilizer interval so that its length cancels the
factor $1+L$. Proposition~\ref{prop:scale-linear-subsequence} provides
$T_n\to\infty$, $A_0<\infty$, and $\eta_0>0$ such that, with fixed
slack in $A_0$, the cutoff $r_T\le A_0T_n$ has lower limiting mass at
least $\eta_0$ times the squared norm of each fixed input, uniformly
for stabilizer lengths in $[T_n/2,T_n]$. The assertion passes to the
constant amplification by tensor density. Since $1_C$ commutes with
the constant-amplified $U_a$, the inputs here are the fixed vectors
$1_CVv_{\gamma j}$. Restrict
\eqref{eq:scale-stabilizer-length-bound} to that stabilizer segment,
set $L=A_0T_n$, and divide by $T_n$. Its Haar measure is comparable to
$T_n$, with the $M$ factor compact. Fatou's lemma, first for finite sums
of $(\gamma,j)$ and then by monotone exhaustion, yields
\begin{equation}
 \sum_{\gamma,j}\langle v_{\gamma j},W_Cv_{\gamma j}\rangle
 \le C'_C\delta^{-Q_T}.
 \label{eq:scale-weighted-frame-bound}
\end{equation}
The constant is independent of the density $f$ satisfying
\eqref{eq:scale-fibre-density}.

Undo induction by lifting at $g=d_0(o')\gamma$ with operator
conjugation by $\pi_0(\gamma)^*$.  Since $W_C$ commutes with $U_S$,
its uninduced field $W_C^0$ depends only on the complementary
coordinate $g'\in\widehat G_S$.  Since the pure-vertex family is
fixed by complementary source factors, its uninduced effect depends
only on the inverse-translated pair.  These two operators commute.
Unfolding the sum over $\gamma$ in
\eqref{eq:scale-weighted-frame-bound} therefore integrates
\[
 \sum_j\langle b_j,
 W_C^0(g')p_\delta^0(x,y)b_j\rangle
\]
against the density induced by $f$ on
$J_0\backslash G\cong B_S^{(2)}\times\widehat G_S$.
Every compactly supported quotient density between zero and one
can be obtained by choosing, on finitely many local sections, a
nonnegative compactly supported fibre density of integral one and
using a partition of unity.  Exhausting the quotient and applying
Tonelli's theorem gives
\begin{equation}
 \int_{\widehat G_S}\int_{B_S^{(2)}}
 \sum_j\langle b_j,W_C^0(g')p_\delta^0(x,y)b_j\rangle
 \,dm(x,y)\,dg'
 \le C'_C\delta^{-Q_T}.
 \label{eq:scale-unfolded-pair-bound}
\end{equation}
The same argument applies to any fixed frame translate
$\pi_0(\gamma_0)W_{\lambda_0}b_j$: its full orbit is again a
Parseval family, up to irrelevant scalar phases.

We spell out why this gives a dense supply of constant inputs for
\eqref{eq:input-10}.  Restrict $o'$ to a measurable set on which
$d_0(o')$ lies in a compact subset $K$ of $G$.  Haar measure on the
section image is restricted Haar, up to the fixed covolume
normalization.  Projection of its restriction to $K$ onto
$\widehat G_S$ is dominated by a finite constant times Haar measure:
the fibre in $S$ lies in a fixed compact set of finite Haar volume.
For a bounded scalar function $a$ supported on this set, put
\[
 \xi(o')=a(o')W_C(o')^{1/2}
             \pi_0(\gamma_0)W_{\lambda_0}b_j.
\]
Commutation with the angular effects moves $W_C^{1/2}$ into the
positive pairing.  Changing the pair variable by the $S$ coordinate
of $d_0(o')$ preserves $dm$.  The preceding Haar domination and
\eqref{eq:scale-unfolded-pair-bound} now prove
\eqref{eq:input-10} for this fixed $\xi$.  Finite sums are
handled by Cauchy--Schwarz.  Such vectors span a dense space:
compact section restrictions exhaust $Y$, the frame translates
span $H$, bounded simple functions are dense in the base variable,
and $W_C^{1/2}\to1$ strongly along a compact exhaustion.

Finally, let $\mathscr E$ be the maximal linear domain of vectors
for which \eqref{eq:input-10} holds for every fixed finite set
of mask controls of sufficiently high order.  The dense supply
just obtained is contained in this domain.  A bounded operator
commuting with all $P_x(f)$ commutes with $p_\delta$ and preserves
the bound, with constant multiplied by its squared norm.  A fixed
source transformation preserves $dm$ and reindexes the pair.  A
fixed target transformation changes scale, support, separation,
and normalized derivatives by fixed constants because it is a
smooth conformal map of the compact horizontal boundary.  Thus it
maps a controlled family into another permitted controlled family.
These observations prove all the asserted domain invariances.
\end{proof}

\subsection{Horizontal differentiation and an injective height}

We retain the matched isolated factors $S,T$ and their boundaries
$B_S,B_T$.  Their homogeneous dimensions are denoted by $Q_S,Q_T$.  The Carnot distances are
$d_S,d_T$, and $dx$ is compact-rotation probability on $B_S$.
For $R=S,T$, let $\mathcal V_R\subset TB_R$ be the horizontal bundle,
with its compact horizontal metric.  Its iterated brackets span $TB_R$.
The action is horizontally conformal; write
$\operatorname{dil}(a,x)>0$ for its horizontal dilation.  Thus
\begin{equation}\label{eq:scale-dilation-cocycle}
 \operatorname{dil}(ab,x)
 =\operatorname{dil}(a,bx)\operatorname{dil}(b,x).
\end{equation}
All the isolated factors in the present setting have $Q_T>2$.

Let $\mathscr E\subset\mathcal H$ be the maximal invariant linear domain
defined by the estimates in Proposition~\ref{prop:scale-pair-estimate}; that
Proposition proves its density.  In particular,
\eqref{eq:input-10} holds for every $\xi\in\mathscr E$, uniformly
over each fixed collection of normalized mask controls.  Constants
below may depend on $\xi$ and these controls.  The domain is invariant
under the source and measured symmetries and under every bounded
operator commuting with all the $P_x$.  We use the strong continuity,
sharpness, covariance, and nonconstancy of the pure-vertex family
$x\mapsto P_x$ established previously.

\begin{definition}\label{def:scale-ordered-kernel}
For $F\in C^\infty(B_T\times B_T)$, its ordered evaluation is
\[
 \mathcal K_{x,y}(F)
 =\iint F(w,v)\,P_y(dw)P_x(dv).
\]
More explicitly, on finitely many coordinate rectangles expand $F$ as
an absolutely summable series of products of smooth functions; replace
$g(w)f(v)$ by $P_y(g)P_x(f)$ and sum in operator norm.  Smooth Fourier
expansion with sufficiently many derivatives gives absolute
summability, with the sum bounded by a fixed sufficiently high smooth
seminorm of $F$.  Approximation in that seminorm by finite separated
smooth sums proves independence of the coordinate expansion.

A smooth kernel has horizontal order at least one if it vanishes on
the diagonal.  It has horizontal order at least two if, in addition,
its derivative in the second variable vanishes there on
$\mathcal V_T$.  These conditions are invariant under the coordinate
changes used below.
\end{definition}

The next estimate uses $\mathcal K_{x,y}$ in its displayed order.
No commutation between $P_x$ and $P_y$ is needed.

\begin{lemma}\label{lem:scale-ordered-differences}
Fix a compactly contained source horospherical chart, a bounded
nonnegative source cutoff $\chi$, and a smooth nonnegative function
$\vartheta$ of compact support on its nilpotent coordinate group.  Set
$y=x\delta_\epsilon w$ in that chart, where $\delta_\epsilon$ is Carnot
dilation, and use the measure
$d\nu_\epsilon=\chi(x)\vartheta(w)\,dx\,dw$.  For $\epsilon$ sufficiently
small, the following conclusions hold.
\begin{enumerate}[label=\textup{(\roman*)}]
\item If $F$ has horizontal order at least one, then, for each
      $\xi\in\mathscr E$,
      \[
       \|\mathcal K_{x,y}(F)\xi\|_{L^2(d\nu_\epsilon;\mathcal H)}
       \le C_{F,\xi}\epsilon^{Q_S/Q_T}.
      \]
\item If $Q_S=Q_T=Q>2$ and $F$ has horizontal order at least two,
      then the same norm is $o(\epsilon)$.  More precisely it is bounded
      by a constant times
      \[
      \begin{cases}
      \epsilon^{Q/2},&2<Q<4,\\
      \epsilon^2(1+|\log\epsilon|),&Q=4,\\
      \epsilon^2,&Q>4.
      \end{cases}
      \]
\end{enumerate}
The constants are uniform when a sufficiently large fixed finite
collection of the indicated smooth seminorms of $F$ is bounded.
\end{lemma}

\begin{proof}
Under the change from $w$ to $y$, the density of $\nu_\epsilon$ with
respect to $dx\,dy$ is at most $C\epsilon^{-Q_S}$ and is supported where
$d_S(x,y)\le C\epsilon$.  Since the invariant pair measure has density
comparable to $d_S(x,y)^{-2Q_S}$, the pullback of any separated target
mask effect satisfies
\begin{equation}\label{eq:scale-averaged-mask-bound}
 \int\langle\xi,p_\delta(x,y)\xi\rangle\,d\nu_\epsilon
 \le C_\xi\epsilon^{Q_S}\delta^{-Q_T}.
\end{equation}
Its bound by a constant times $\|\xi\|^2$ also follows from the uniform
operator bound for $p_\delta$ in Definition~\ref{def:scale-mask-controls}.

Here are the details needed to apply this estimate to a kernel.
Choose a smooth dyadic decomposition of the complement of the target
diagonal.  The piece at scale $\delta$ is supported where
$c\delta\le d_T(w,v)\le C\delta$.  Cover each boundary by cells of
radius $c'\delta$ with $c'$ small relative to $c$.  The pairs of cells
meeting this support have bounded degree: each cell is paired with at
most a fixed number of other cells.  Their enlarged supports are
separated, and their centers have uniformly bounded density at scale
$\delta$.

Suppose $F$ has horizontal order at least $m$, where $m=1$ or $2$.
On each such product of cells, every fixed finite collection of
normalized derivatives of the kernel piece is bounded by $C\delta^m$.
For $m=1$, this is Taylor expansion at the diagonal.  For $m=2$, the
constant and horizontal linear terms vanish; the remaining horizontal
quadratic terms and the nonhorizontal linear terms both have Carnot
degree at least two.  Taylor expansion in the rescaled coordinates
also gives the same estimate after each of the specified normalized
derivatives.  Compactness makes the bounds uniform in the cell centers.

Expand the rescaled kernel in a Fourier series on slightly larger
coordinate boxes.  Integration by parts in both variables, to an
order larger than the box dimensions and the required mask orders,
gives a summable family of separated terms
\[
 \delta^m c_{A,D;k,l}\,
       g_{D,k}(w)f_{A,l}(v).
\]
The coefficients decrease faster than the polynomial growth of the
normalized mask seminorms in $k,l$.  For a fixed pair of frequencies,
the leftmost supports have bounded overlap.  The spectral calculus
for $P_y$, followed by Cauchy--Schwarz for the bounded number of terms
at a cell, therefore gives
\[
 \left\|\sum_{A,D}c_{A,D;k,l}
       P_y(g_{D,k})P_x(f_{A,l})\xi\right\|^2
 \le C_{k,l}\sum_{A,D}
       \|P_y(g_{D,k})P_x(f_{A,l})\xi\|^2.
\]
Coefficients can be absorbed into the masks; duplicating an index
when necessary changes density and overlap by only the bounded
degree.  Thus \eqref{eq:scale-averaged-mask-bound}, followed by the
summation over frequencies, bounds this annular piece in root mean
square by
\begin{equation}\label{eq:scale-annular-kernel-bound}
 C_{F,\xi}\delta^m
       \min\{1,\epsilon^{Q_S/2}\delta^{-Q_T/2}\}.
\end{equation}
This computation uses only one PVM for the leftmost overlap and the
ordered positive effect for the remaining factors.

There is also an estimate $C_F\delta_0^m\|\xi\|$ for the entire kernel
piece supported where $d_T(w,v)\le C\delta_0$.  Indeed repeat the cell
expansion at scale $\delta_0$; only boundedly many cells meet a given
cell, and the constant mask bound applies without separation.

Put $a=Q_S/Q_T$ and choose a dyadic $\delta_0$ comparable to
$\epsilon^a$.  For $m=1$, the triangle inequality and
\eqref{eq:scale-annular-kernel-bound} give
\[
 C_{F,\xi}\left(\delta_0+
       \epsilon^{Q_S/2}\sum_{\delta\ge\delta_0}
                         \delta^{1-Q_T/2}\right)
 \le C'_{F,\xi}\epsilon^a,
\]
because $Q_T>2$.  When $Q_S=Q_T=Q$ and $m=2$, take
$\delta_0\asymp\epsilon$ instead.  Summing
$\epsilon^{Q/2}\delta^{2-Q/2}$ gives exactly the three rates in the
statement, including the $C\epsilon^2$ diagonal remainder.  Each rate
is $o(\epsilon)$.
\end{proof}

\begin{lemma}\label{lem:scale-equal-dimension}
For every matched isolated pair, $Q_S=Q_T$.  Write their common value
as $Q$.  If $X$ is a smooth horizontal source field,
$f\in C^\infty(B_T)$, and $\xi\in\mathscr E$, then the distributional
derivative
\[
 D_X(f)\xi=X\bigl(x\mapsto P_x(f)\xi\bigr)
\]
belongs to $L^2(B_S,dx;\mathcal H)$.  Differentiation is linear in
$f,\xi$, and is linear over smooth source functions in $X$.
\end{lemma}

\begin{proof}
For $F(w,v)=f(w)-f(v)$, ordered evaluation is
$P_y(f)-P_x(f)$.  To extract a horizontal derivative, use a smooth
compactly supported weight $\psi_i(w)$ whose horizontal first moments
are the $i$th coordinate vector.  Its absolute value is dominated by
a fixed smooth nonnegative weight of the kind used in
Lemma~\ref{lem:scale-ordered-differences}.  The expressions
\[
 \frac1\epsilon\int\psi_i(w)
            (P_{x\delta_\epsilon w}(f)-P_x(f))\xi\,dw
\]
are consequently bounded in local $L^2$ by
$C\epsilon^{Q_S/Q_T-1}$.  For a smooth compactly supported test field,
change variables from $x$ to $x\delta_\epsilon w$ in its adjoint
pairing.  First-order Taylor expansion of the test field and the
smooth measure density shows that the limit is the distributional
horizontal derivative: the adjoint test is
$-X\varphi-(\operatorname{div}_{dx}X)\varphi$.

If $Q_S/Q_T>1$, all these derivatives vanish.  Brackets of horizontal
fields then annihilate the same distribution, because
$[X,Y]u=X(Yu)-Y(Xu)$.  Bracket generation implies that every coordinate
derivative vanishes, so $P_x(f)\xi$ is constant on the connected
boundary.  Density of $\mathscr E$ and strong continuity make $P_x$
constant, a contradiction.  Hence $Q_S\le Q_T$.  The matching is a
permutation of the finite collection of isolated types; going around
each cycle forces equality throughout that cycle.

With equality, the displayed difference quotients are bounded in
$L^2$.  Weak compactness and their identified distributional limit
give $D_X(f)\xi\in L^2$ in each chart.  A finite cover and a smooth
partition of unity give the global statement.  The linearity
assertions follow directly from distributional differentiation.
\end{proof}

We next specify the operator domain.  This is necessary because an
individual vector $P_x(g)\xi$ need not belong to the constant-vector
estimate domain $\mathscr E$.

\begin{lemma}\label{lem:scale-module-extension}
Set $\mathcal K=L^2(B_S,dx;\mathcal H)$ and
\[
 \mathcal C=\operatorname{span}\{x\mapsto a(x)P_x(g)\xi:
       a\in L^\infty(B_S),\ g\in C^\infty(B_T),\ \xi\in\mathscr E\}.
\]
For real $f$ and real horizontal $X$, the rule
\begin{equation}\label{eq:scale-module-extension-rule}
 T_X(f)\left(\sum_j a_j P_{\bullet}(g_j)\xi_j\right)
 =\sum_j a_j P_{\bullet}(g_j)D_X(f)\xi_j
\end{equation}
defines a densely defined symmetric operator on $\mathcal C$.
Furthermore, if real smooth $g_j,f_j$ satisfy
$\sum_j g_j\,d_{\rm hor}f_j=0$, then
\begin{equation}\label{eq:scale-horizontal-cancellation}
 \sum_j P_{\bullet}(g_j)T_X(f_j)=0
 \quad\hbox{on }\mathcal C.
\end{equation}
All the operators in \eqref{eq:scale-module-extension-rule} commute on
$\mathcal C$ with scalar parameter multipliers and smooth angular
multipliers.  Every real smooth angularly weighted component
$P_{\bullet}(g)T_X(f)$ is symmetric on $\mathcal C$.
\end{lemma}

\begin{proof}
Write $\Delta_{x,y}f=P_y(f)-P_x(f)$.  Sharpness gives the exact identity
\[
 (\Delta_{x,y}f)P_x(g)-P_y(g)\Delta_{x,y}f
 =\mathcal K_{x,y}\bigl((f(w)-f(v))(g(v)-g(w))\bigr).
\]
The right side has horizontal order two.  Its norm on any
$\xi\in\mathscr E$, divided by $\epsilon$, tends to zero in the
averages used to recover a horizontal derivative.

Fix $\xi,\eta\in\mathscr E$ and real $g$.  In the pairing of the left
term with $\eta$, move $\Delta_{x,y}f$ to $\eta$ by adjunction.
For the other term, keep the bounded quotient applied to $\xi$ and
move $P_y(g)$ to the test vector.  Strong continuity permits replacing
$P_y(g)$ there by $P_x(g)$: its difference on a fixed test vector tends
to zero in $L^2$ by dominated convergence, whereas the averaged
quotients stay bounded in $L^2$.  The weak derivative limit therefore
gives
\begin{equation}\label{eq:scale-weighted-matrix-symmetry}
 \langle P_x(g)D_X(f)(x)\xi,\eta\rangle
 =\langle\xi,P_x(g)D_X(f)(x)\eta\rangle
 \quad\hbox{for almost every }x.
\end{equation}
One obtains the pointwise assertion by first inserting arbitrary
bounded scalar parameter tests.  Complex weights follow by linearity
and adjunction.

Apply \eqref{eq:scale-weighted-matrix-symmetry} with products of angular
tests.  If $u=\sum_j a_jP_{\bullet}(g_j)\xi_j$ and
$v=\sum_k b_kP_{\bullet}(h_k)\eta_k$, the formal right side of
\eqref{eq:scale-module-extension-rule} satisfies
\[
 \langle T_X(f)u,v\rangle_{\mathcal K}
 =\langle u,T_X(f)v\rangle_{\mathcal K}.
\]
In particular, if an expression for $u$ represents zero, its formal
image is orthogonal to every $v\in\mathcal C$.  This space is dense,
since it includes bounded simple parameter fields with values in
$\mathscr E$.  The image is therefore zero.  This proves both
well-definedness and symmetry, and hence closability.  The identical
calculation with one more real angular weight proves the final
symmetry assertion.

For cancellation, consider
\[
 F(w,v)=\sum_j g_j(w)(f_j(w)-f_j(v)).
\]
It vanishes on the diagonal, and its second-variable horizontal
derivative there is $-\sum_jg_jd_{\rm hor}f_j=0$.
Lemma~\ref{lem:scale-ordered-differences} gives an
$o(\epsilon)$ bound.  Taking the same weak derivative limit and
replacing left multipliers at $y$ by those at $x$ gives
$\sum_jP_x(g_j)D_X(f_j)(x)\xi=0$ for every $\xi\in\mathscr E$.
Commuting the additional multipliers in
\eqref{eq:scale-module-extension-rule} proves
\eqref{eq:scale-horizontal-cancellation} on $\mathcal C$.  The asserted
commutation on this domain also follows from that rule.
\end{proof}

For a real smooth horizontal target covector field
$\theta=\sum_jg_jd_{\rm hor}f_j$, define
\[
 T_X(\theta)=\sum_jP_{\bullet}(g_j)T_X(f_j)
 \quad\hbox{on }\mathcal C.
\]
Equation~\eqref{eq:scale-horizontal-cancellation} makes this independent
of the chosen expression.  Such expressions exist globally: finitely
many smooth coordinate functions have horizontal differentials
spanning the bundle, and their Gram operator has a smooth inverse on
that bundle.

\begin{lemma}\label{lem:scale-closed-covariance}
There is a positive self-adjoint measurable field $l_x$ on
$\mathcal H$, affiliated to $\mathcal A_1$, commuting strongly with
$P_x$, and invariant under the complementary source factors, such
that, for every fixed source element $a\in S$ and measured lattice
element $h$,
\begin{equation}\label{eq:scale-derivative-modulus-covariance}
 \sigma_a(l_x)=\operatorname{dil}(a,x)l_{ax},\qquad
 \beta_h(l_x)=\operatorname{dil}(h^{-1},P_x)l_x
\end{equation}
for almost every $x$.  The field is not identically zero.
Here the target dilation is evaluated by the spectral calculus of
$P_x$, using the $T$-coordinate of $h$.
\end{lemma}

\begin{proof}
Choose finite smooth Parseval horizontal source fields
$(X_\alpha)_{\alpha=1}^{n_S}$ and finite smooth Parseval horizontal
target covectors $(\theta_r)_{r=1}^{n_T}$.
For completeness, start with finitely many smooth spanning sections.
Their positive Gram operator is invertible on the horizontal bundle;
multiplication by its inverse square root gives the required Parseval
system.  Smooth coordinate differentials provide the target spanning
sections, so each $\theta_r$ has an expression of the form used above.
Define the finite column
\[
 D:\mathcal C\longrightarrow\mathcal K^{n_Sn_T},\qquad
 Du=(T_{X_\alpha}(\theta_r)u)_{\alpha,r}.
\]
Every component is symmetric on $\mathcal C$ by
Lemma~\ref{lem:scale-module-extension}.  If $u_n\to0$ and
$Du_n\to v$, testing each component against $\mathcal C$ gives $v=0$.
Thus $D$ is closable.  Let $\overline D$ be its closure and put
$l=|\overline D|=(\overline D^*\overline D)^{1/2}$ on $\mathcal K$.

Multiplication by a scalar parameter unitary preserves $\mathcal C$
and commutes with $D$, acting diagonally on the output.  The same is
true for $P_{\bullet}(e^{it g})$, with $g$ real smooth.  Their inverses
have the same property.  These identities extend to the closed graph,
and consequently to the closed quadratic form
$u\mapsto\|\overline D u\|^2$.  Therefore $l$ strongly commutes with
the parameter multiplication algebra and the angular spectral algebra.
In particular $l=\int^{\oplus}l_x\,dx$ is decomposable and $l_x$
strongly commutes with $P_x$ almost everywhere.

If $V$ is a unitary commuting with every $P_x$, invariance of
$\mathscr E$ and distributional differentiation give
$D_X(f)V\xi=VD_X(f)\xi$.  This identity extends to $\mathcal C$ and
then to the closed graph.  Applying it to the base unitaries and the
unmeasured group actions proves affiliation to $\mathcal A_1$;
applying it to the complementary source factors proves their
invariance.  Separability permits a common determining set for these
almost-everywhere commutation assertions.

We give the norm argument for covariance to keep track of closures.
The Parseval systems realize isometric embeddings of the horizontal
bundles into trivial finite-dimensional bundles.  Their Gram matrices
are the corresponding orthogonal projections.  Linearity in $X$ and
\eqref{eq:scale-horizontal-cancellation} say exactly that $Du$ belongs
to the range of these source and spectral target projections.  Thus
an isometry between horizontal bundles preserves the sum of squares
of the column entries, even if the Parseval systems are redundant.
In particular, any two finite Parseval systems give the same squared
form on the same intrinsic module core.  Their closed squared forms,
and hence their operators $l$, agree.  The maximal estimate domain
$\mathscr E$ is also unchanged by replacing the controlled Carnot
atlas: a finite refinement and the smooth changes of rescaled
coordinates transfer every mask control with a fixed change of its
constants, in both directions.  Thus the construction depends only on
the PVM family and the chosen compact horizontal metrics.  The cell
expansions used to prove the estimates, and the localizing isometry
used to prove density of $\mathscr E$, do not enter its definition.

Differentiate the identity $\sigma_aP_x=P_{ax}$ on constant-domain
vectors.  The chain rule sends a source vector $X$ at $x$ to
$da_xX$ at $ax$.  This map is the product of
$\operatorname{dil}(a,x)$ and a horizontal isometry.  Similarly,
differentiating $\beta_hP_x(f)=P_x(f\circ h^{-1})$ replaces a target
covector by its pullback under $h^{-1}$.  Horizontal cancellation
justifies writing this pullback in the chosen target coframe.  Its
norm factor, evaluated spectrally, is
$\operatorname{dil}(h^{-1},P_x)$; its other factor is a horizontal
isometry.  The preceding projection identities justify dropping those
isometries when taking the sum of squared column norms.

These identities hold on the full core $\mathcal C$.  Indeed source
transport reindexes its scalar coefficients, transforms
$\mathscr E$ into itself, and transports $P_x(g)$ by covariance;
measured transport replaces $g$ by $g\circ h^{-1}$ and also preserves
$\mathscr E$.  On $\mathcal K$, source change of variables includes
the square root of the smooth Radon--Nikodym derivative of $dx$.
This is an allowed scalar multiplier and commutes with $D$, so it
contributes no additional derivative.  For each fixed group element,
all dilation factors and inverses are bounded on the compact
boundaries.  Source and target transforms therefore give invertible
maps of the graph norms.  The identities extend to the closed
quadratic forms.  The uniqueness of the positive operator associated
to such a form proves
\eqref{eq:scale-derivative-modulus-covariance}; spectral target
dilations commute strongly with $l$, as already proved.

Finally, if $l=0$, every column component vanishes on $\mathcal C$.
Spanning by the Parseval fields and covectors, followed by horizontal
cancellation, gives $D_X(f)\xi=0$ for every smooth $f$, horizontal $X$,
and $\xi\in\mathscr E$.  The bracket argument in
Lemma~\ref{lem:scale-equal-dimension} would make $P_x$
constant.  Hence $l\ne0$.
\end{proof}

\begin{proposition}[Height at an isolated type]\label{prop:scale-isolated-height}
For a matched isolated pair the homogeneous dimensions agree,
$Q_S=Q_T=Q$.  There is a family of positive injective self-adjoint
operators $b_x$, affiliated to $\mathcal A_1$, for every $x\in B_S$,
which strongly commute with $P_x$ and are invariant under the
complementary source factors.  They satisfy the exact identities
\begin{equation}\label{eq:scale-isolated-height-covariance}
 \sigma_a(b_x)=\operatorname{dil}(a,x)^{-Q}b_{ax},\qquad
 \beta_h(b_x)=\operatorname{dil}(h^{-1},P_x)^{-Q}b_x.
\end{equation}
These are the inverse angular-Jacobian laws.  The same construction
passes to the normal copies used in the amplified models.
\end{proposition}

\begin{proof}
We first replace the field in
Lemma~\ref{lem:scale-closed-covariance} by an everywhere covariant
version.  Fix $x_0\in B_S$ and let $J$ be its parabolic stabilizer.
Pull the field back along $g\mapsto gx_0$ and set
\[
 F(g)=\operatorname{dil}(g,x_0)\sigma_g^{-1}(l_{gx_0}).
\]
This is a measurable field of positive self-adjoint operators.  One
can test measurability through their bounded resolvents, or their
closed graph projections.  Haar disintegration over $S/J$ and
\eqref{eq:scale-derivative-modulus-covariance} give
$F(ag)=F(g)$ for almost every $g$, for each fixed $a\in S$.
Fubini, followed by the nonsingular change of variables
$(a,g)\mapsto(ag,g)$, shows that $F(g_1)=F(g_2)$ for almost every pair.
Thus $F$ has a single essential value $L_0$.  A countable determining
family of bounded matrix coefficients of graph projections suffices
for this argument.

For each $k\in J$ the definition and
\eqref{eq:scale-dilation-cocycle} give
\[
 F(gk)=\operatorname{dil}(k,x_0)\sigma_k^{-1}F(g).
\]
Consequently $\sigma_kL_0=\operatorname{dil}(k,x_0)L_0$.  We may
therefore define, independently of the representative $g$,
\[
 l_{gx_0}=\operatorname{dil}(g,x_0)^{-1}\sigma_gL_0.
\]
This agrees with the old field almost everywhere and satisfies its
source covariance for every $g,x$.  Its other properties hold as well.
To see the target law explicitly, commute $\beta_h$ past $\sigma_g$
in the pullback.  Since $\sigma_gP_{x_0}=P_{gx_0}$, its spectral
dilation becomes the fixed operator
$\operatorname{dil}(h^{-1},P_{x_0})$ on $F(g)$.  The essential value
$L_0$ therefore satisfies the target law at $x_0$, and transport proves
it everywhere.  Angular commutation, affiliation, and complementary
source invariance pass to $L_0$ and transport in the same way.

Let $e_0=\supp L_0$.  The target dilation is bounded, positive,
invertible, and commutes with $L_0$.  Taking supports in its covariance
law therefore gives $\beta_h(e_0)=e_0$ for all $h$, so $e_0\in\mathcal N$.
Source stabilizer covariance gives $\sigma_k(e_0)=e_0$ for $k\in J$.
In addition, $e_0$ is fixed by every complementary source factor.
Choose a split sequence in $J$ escaping in $S$, and multiply it by
sequences escaping in each complementary factor.  Every resulting
element fixes $e_0$, while the all-factor mixing established in Section~\ref{sec:setup}
makes its translates converge weakly to $\tau(e_0)1$.  Hence
$e_0=\tau(e_0)1$.  It is nonzero because the derivative column was
nonzero; as a projection it must be $1$.  This use of complementary
source invariance is essential to the application of mixing.

All transported $l_x$ thus have zero kernel.  Define
$b_x=l_x^{-Q}$ by the spectral calculus.  It is densely defined,
positive, self-adjoint, injective, and affiliated to $\mathcal A_1$.
Taking inverse powers in the two covariance laws, using strong
commutation with the spectral dilation, proves
\eqref{eq:scale-isolated-height-covariance}.  The compact horizontal
Jacobian is the $Q$th power of the dilation, so these are precisely
the asserted inverse-Jacobian laws.

We verify the normal-copy assertion at the level of domains.  For an
ampliation $\mathcal H\otimes\mathcal L$, fix an orthonormal basis
$(e_n)$ of the separable multiplicity space.  If
$\Xi=\sum_n\xi_n\otimes e_n$ belongs to its maximal estimate domain,
positivity of every pair effect implies $\xi_n\in\mathscr E$ for each
$n$: each coordinate energy is bounded by the total energy.  Conversely,
every finite sum with $\xi_n\in\mathscr E$ belongs to the amplified
estimate domain.  The finite-coordinate projections commute with the
PVMs, preserve that domain, and commute with its weak derivatives.
For a generator $aP_{\bullet}(g)\Xi$ of the amplified module core,
those projections therefore converge in both the input norm and the
derivative-column norm.  Its projected vectors lie in
$\mathcal C\odot\operatorname{span}\{e_n\}$, and the same holds for finite sums of
generators.  Thus this algebraic tensor core is graph-dense in the
amplified core.  On it the derivative column is exactly $D\otimes1$.
Their closures agree, and the spectral calculus gives
$l^{\rm amp}=l\otimes1$ and $b^{\rm amp}=b\otimes1$.  Unitary
identification with a normal copy preserves these identities.  In
particular no mixing assertion on an enlarged commutant is required.
\end{proof}


The isolated branch now supplies its own-label height calculus:
$b_x$ is positive and injective, commutes strongly with $P_x$, and obeys
\eqref{eq:scale-isolated-height-covariance} in every required normal copy.
Section~\ref{sec:heights} uses these identities to transport imaginary
powers through the Fourier links and establish the centered length laws.
The closed derivative column above is the construction of this height;
the common graph domain used to compare heights at different labels is
constructed later from the Haar-height isometry in
Section~\ref{sec:locality}.
